% !TeX root = ../tesis.tex
%%%%%%%%%%%%%%%%%%%%%%%%%%%%%%%%%%%%%%%%%%%%%%%%%%%%%%%%%%%%%%%%%%%
% SECCIÓN 5.2 — CONSTRUCCIÓN DEL GRAFO TOPOLÓGICO
% Issue #24 [C-7]
%%%%%%%%%%%%%%%%%%%%%%%%%%%%%%%%%%%%%%%%%%%%%%%%%%%%%%%%%%%%%%%%%%%
\section{Construcción del Grafo Topológico}
\label{sec:grafo-topologico}

A partir de los datos geoespaciales procesados por el \textbf{Apímetro}, el
\textbf{Modelo VFT} construyó un grafo dirigido y ponderado que representa la
red de transporte masivo de la \zmvm{} como objeto de análisis estructural
(véanse Capítulo~4 y Anexo~B para el detalle de implementación). El proceso de
integración topológica —mediante un árbol \texttt{KDTree} con tolerancia de
85~m, corrida del 18 de abril de 2026— redujo las 22{,}766 entidades espaciales
del dataset de origen a \textbf{11{,}115 nodos canónicos} conectados por
\textbf{45{,}679 aristas}, de las cuales 25{,}197 corresponden a segmentos de
servicio real y 20{,}482 a aristas de transbordo peatonal generadas
algorítmicamente \citep{cabrera2026VFTModel}.

\begin{figure}[H]
    \centering
    % Diagrama: Pipeline de procesamiento interno del Modelo VFT (4 etapas)
% Usado en: 5-Resultados/5-1-Introduccion-Resultados.tex  (fig:arq-capas)
% Reemplaza: 3 subfiguras arq_capa1/2/3.png con captions incorrectos
\begin{tikzpicture}[
    node distance = 0.6cm and 0.5cm,
    etapa/.style = {
        rectangle, draw=unamAzul, fill=unamAzul!8,
        text width=2.8cm, align=center,
        minimum height=1.5cm, rounded corners=4pt,
        font=\small\bfseries
    },
    dato/.style = {
        font=\scriptsize\itshape, color=unamAzul!75,
        align=center, text width=2.8cm
    },
    cifra/.style = {
        font=\footnotesize\bfseries, color=unamAzul,
        align=center, text width=2.8cm
    },
    flecha/.style = {-Stealth, thick, color=unamAzul!65}
]

%% — Etapas —
\node[etapa] (ing) {Ingesta\\GTFS};
\node[etapa, right=of ing]  (snap) {Snapping\\Lógico};
\node[etapa, right=of snap] (grf)  {Construcción\\del Grafo};
\node[etapa, right=of grf]  (imp)  {Ponderación\\Impedancia};

%% — Cifras sobre cada etapa —
\node[cifra, above=0.18cm of ing]  {22{,}766 entidades};
\node[cifra, above=0.18cm of snap] {11{,}115 nodos};
\node[cifra, above=0.18cm of grf]  {45{,}679 aristas};
\node[cifra, above=0.18cm of imp]  {grafo $G$ dirigido};

%% — Detalles bajo cada etapa —
\node[dato, below=0.18cm of ing]  {Apimetro REST\\GeoJSON};
\node[dato, below=0.18cm of snap] {KDTree · $d \leq 50$~m\\eliminación duplicados};
\node[dato, below=0.18cm of grf]  {NetworkX\\25{,}197 servicio\\20{,}482 transbordo};
\node[dato, below=0.18cm of imp]  {Motor VFT\\pesos en minutos};

%% — Flechas —
\draw[flecha] (ing)  -- (snap);
\draw[flecha] (snap) -- (grf);
\draw[flecha] (grf)  -- (imp);

%% — Nota conceptual: abstracción Red Física → Grafo Matemático —
\node[font=\scriptsize\color{gray!70}, anchor=north]
    at ($(ing.south)!0.5!(imp.south) + (0,-1.15cm)$)
    {Red Física $\longrightarrow$ Grafo Matemático $\mathcal{G}(\mathcal{V},\,\mathcal{E},\,w)$};

\end{tikzpicture}

    \caption[Pipeline de procesamiento del Modelo VFT]{Las cuatro etapas que
    transforman los datos brutos de \textbf{Apímetro} en el grafo topológico
    operativo. El \termino{Snapping Lógico Geográfico} mediante árbol
    \texttt{KDTree} (tolerancia: 85~m) reduce las 22{,}766 entidades originales
    a 11{,}115 nodos canónicos; la ponderación final asigna impedancias
    temporales a las 45{,}679 aristas del grafo dirigido
    $\mathcal{G}(\mathcal{V},\mathcal{E},w)$.}
    \label{fig:arq-capas}
    \fuentefigura{Elaboración propia con \texttt{VFTModel}
    \citep{cabrera2026VFTModel}.}
\end{figure}

\begin{itemize}
    \item \textbf{11{,}115 nodos topológicos} — estaciones operativas tras
    eliminar nodos fantasma y redundancias geométricas.
    \item \textbf{45{,}679 aristas}: 25{,}197 segmentos de servicio real y
    20{,}482 aristas de transbordo peatonal generadas por el
    \termino{Snapping Lógico Geográfico}.
    \item \textbf{11 sistemas de transporte} representados: \textsc{Metro},
    Metrobús, Trolebús, Tren Ligero, Tren Suburbano, Cablebús, Mexicable,
    Mexibús, \textsc{rtp}, Corredores Concesionados y Pumabús.
    \item Distribución modal: \textbf{12\,\% infraestructura estructurada} con
    derecho de vía exclusivo o confinado (\textsc{Metro}, Metrobús, Tren Ligero,
    Cablebús, Tren Suburbano) y \textbf{88\,\% transporte de superficie} en
    vía compartida o mixta (\textsc{rtp}, Corredores Concesionados, Trolebús,
    Mexibús, Mexicable, Pumabús).
    \item \textbf{Componente gigante}: 10{,}561 nodos (95\,\% del total),
    sobre la cual se ejecutan los algoritmos de Fase~3.
\end{itemize}

El Cuadro~\ref{tab:aristas-por-sistema} detalla la distribución de aristas de
servicio real por sistema, el tipo de derecho de vía asignado y el coeficiente
de fricción vial correspondiente.

\begin{table}[H]
    \centering
    \caption[Distribución de aristas de servicio por sistema de transporte]{
    Distribución de aristas de servicio real por sistema de transporte en el
    grafo de la \zmvm{}. El coeficiente $CF$ refleja el tipo de derecho de vía.
    Corrida: 2026-04-18.}
    \label{tab:aristas-por-sistema}
    \begin{tabular}{|l|r|l|c|}
        \hline
        \textbf{Sistema} & \textbf{Aristas} & \textbf{Derecho de vía} &
        \textbf{$CF$} \\
        \hline
        \textsc{rtp}                  & 11{,}041 & Mixto       & 1.759 \\
        Corredores Concesionados (CC) & 10{,}603 & Mixto       & 1.759 \\
        Trolebús\textsuperscript{†}   &  1{,}506 & Compartido  & 1.380 \\
        Metrobús (MB)                 &    792   & Exclusivo   & 1.000 \\
        \textsc{Metro} (STC)          &    528   & Exclusivo   & 1.000 \\
        Mexibús\textsuperscript{†}    &    391   & Confinado   & 1.152 \\
        Pumabús                       &    234   & Mixto       & 1.759 \\
        Tren Ligero (TL)              &     52   & Exclusivo   & 1.000 \\
        Cablebús (CBB)\textsuperscript{‡} &  30  & Confinado   & 1.152 \\
        Mexicable\textsuperscript{†}  &     17   & Confinado   & 1.152 \\
        Tren Suburbano (SUB)          &      3   & Exclusivo   & 1.000 \\
        \hline
        \textbf{Total}                & \textbf{25{,}197} & & \\
        \hline
    \end{tabular}
    \vspace{0.5em}

    \noindent\textsuperscript{†}\,A la fecha de ingesta (2026-04-18), la
    Línea 13 del Trolebús, el Mexicable y el Mexibús no contaban con registros
    publicados en el Portal de Datos Abiertos de la \cdmx{}. Su geometría y
    paradas fueron integradas a partir de datos de
    \termino{OpenStreetMap} \citep{osm2012}.

    \noindent\textsuperscript{‡}\,El sistema Cablebús registra únicamente
    30 aristas, cifra inferior a la esperada para sus tres líneas operativas.
    Esta subrepresentación es una limitación del \textit{backend} del
    \textbf{Apímetro} en la fecha de corrida y no afecta los indicadores de
    los sistemas restantes.

    \fuentefigura{Elaboración propia con \texttt{VFTModel}
    \citep{cabrera2026VFTModel}; datos: \citep{adip_gtfs_2024, osm2012}.}
\end{table}

\begin{figure}[H]
    \centering
    \includegraphics[width=0.92\textwidth]{%
        Figures/Mapas/Qgis_Resources/RedActual/Mapa_Transporte_Publico}
    \caption[Red multimodal de la \zmvm{} por sistema de transporte]{Red
    multimodal de la \zmvm{} organizada por tipo de sistema de transporte,
    elaborada con \textbf{QGIS} a partir de los datos procesados por el
    \textbf{Apímetro}. La versión cartográfica extendida de la red esquemática
    se presenta en la Figura~\ref{fig:anx-mapa-red-esquematica} del
    Anexo~\ref{anx:cartografia-presentacion}.}
    \label{fig:mapa-transporte-publico}
    \fuentefigura{Elaboración propia con \textbf{QGIS};
    datos: \citep{adip_gtfs_2024, osm2012}.}
\end{figure}

El grafo resultante constituye la representación operativa sobre la que se
calculan todos los indicadores presentados en las secciones siguientes. Su
escala —11{,}115 nodos y 45{,}679 aristas distribuidos en once sistemas— hace
de la \zmvm{} uno de los grafos de transporte público más densos documentados
en la literatura latinoamericana \citep{lotero2014vulnerabilidad}, lo que
confiere mayor relevancia a los resultados de los indicadores de cobertura,
centralidad y vulnerabilidad que se exponen a continuación.
